% !PRINTED 389
\weberpage{389}
%<<CONT
bruchentwickelung schliesslich dieselben Schlusszahlen bekommen.
Sind also beide Kettenbrüche periodisch, so müssen sie derselben
Periode angehören.

Hat man für eine gegebene Discriminante $D$ nach den am
Anfang des Paragraphen gegebenen Vorschriften das vollständige
System der reducirten Zahlen entwickelt, so ist es leicht, diese
Zahlen in Perioden zu ordnen und die Theilnenner $\alpha$ des Kettenbruches
 zu finden. Es sei nämlich
\weq{12}{\omega = \frac{\sqrt{D}+b}{2c} = \frac{2a}{\sqrt{D}-b}}
eine von diesen Zahlen und
\weq{13}{\omega_1 = \frac{\sqrt{D}+b_1}{2c_1} = \frac{2a_1}{\sqrt{D}-b_1}}
die ihr in der Periode unmittelbar folgende Zahl, so dass
\weq{14}{\omega = \alpha + \frac{1}{\omega_1}, \quad \omega_1 = \frac{1}{\omega-\alpha}}
ist, wenn $\alpha$ die grösste in $\omega$ enthaltene ganze Zahl bedeutet.
Nach (12) und (14) ist dann
\[
\omega_1 = \frac{2c}{\sqrt{D}+b-2c\alpha},
\]
und die Vergleichung mit (13) giebt
\weq{15}{a_1 = c, \quad b_1 = 2c\alpha - b.}
Wenn umgekehrt zwei der zu $D$ gehörigen reducirten Zahlen
$\{a, b, c\}$, $\{a_1, b_1, c_1\}$ in der durch (15) dargestellten Beziehung
stehen, worin $\alpha$ irgend eine ganze Zahl ist, so folgt die zweite
der ersten in der Periode unmittelbar nach. Denn aus (15)
folgt (14), und daher muss $\alpha$ nach 2. die grösste in $\omega$ enthaltene
Zahl sein.

Die Zahl $\omega_1$ ist also aus $\omega$ durch die Bedingungen (15)
vollständig und eindeutig bestimmt.

7. \gesperrt{Man ordnet also die Zahlen $\{a, b, c\}$ von links
nach rechts in der Weise, dass die letzte Zahl $c$
der vorangehenden zugleich die erste Zahl $a_1$ der
folgenden wird, und dass die Summe der beiden
mittleren Zahlen $b + b_1$ durch $2c$ theilbar ist.
Diese Anordnung ist nur auf eine Art möglich,
und der Quotient $b + b_1 : 2c$ ist die Zahl $\alpha$, die
als Theilnenner im Kettenbruch auftritt.}
